\chapter{背景与动机}
\label{chap:background}

端侧运行时按 collective 语义组织通信计划，底层传输按连接推进发送与反馈，
在网聚合则按通信组和 packet offset 维护共享状态。三者处理同一批数据，却
各有一套完成与资源边界。本章从在网聚合的收益出发（\cref{ssec:p2p-ccl}），
先解释专用聚合状态为何不能直接转化为活跃任务可用的容量
（\cref{ssec:inc-survey}），再考察现有方案如何把组内协调绑定在特定传输
边界上（\cref{ssec:group-consistency}），最后提炼传输层必须满足的上下文
与反馈要求（\cref{ssec:inc-limits}）。

\section{集合通信与在网聚合}
\label{ssec:p2p-ccl}
\label{ssec:inc-mech}

一次集合通信由一组 rank 按共同语义交换数据：\texttt{AllReduce}、
\texttt{Reduce} 和 \texttt{ReduceScatter} 对对应输入执行结合性规约，
\texttt{Broadcast} 和 \texttt{AllGather} 分发一个或多个 rank 的数据，
\texttt{Barrier} 同步参与者。一次调用给出 buffer、数据范围、类型和可选
规约算子，调用完成时各 rank 得到该原语规定的输出。

NCCL~\cite{nccl} 等通信库用 ring、tree 或分层算法实现这些语义；
MSCCL~\cite{msccl}、HiCCL~\cite{hiccl} 和 TACCL~\cite{taccl} 允许进一步
组合发送、接收、规约和拷贝动作。生成的计划规定数据块、参与 rank 和动作
依赖，执行层再把每个动作交给点对点传输。

点对点算法同时决定链路上的数据量。对大小为 $M$ 的 $N$-rank ring
AllReduce，reduce-scatter 和 all-gather 各使每个 rank 发送 $(N-1)M/N$
字节，合计 $2(N-1)M/N$。$N=16$ 时每个 rank 要发送 $1.875M$；仅考虑
100-Gbps 主机出口的串行化，logical goodput 上界就只剩 53.3~Gbps。

\parab{点对点执行以连接为状态单位.}
每条传输分别维护路径、拥塞状态、序号和可靠性，collective 运行时根据完成
事件推进后续动作。这一分工让通信库自由组合计划，也使传输反馈天然对应一个
sender--receiver 关系；同一 collective 内的不同连接同样按各自状态调速和
选路。

\parab{整组关系留在 collective 计划中.}
同一 packet offset 的多份输入共享通信组、聚合算子和完成条件，一个输入的
进度会影响整组数据何时进入下一步。运行时保存这些依赖，网络报文则沿各自
连接传输。一旦网络设备参与 fan-in，它就需要为这些输入选择共同的聚合树和
状态，并从多个分支的到达事件判断整组是否完成——聚合与恢复的边界因而大于
任何一条单独连接。

在网聚合让交换机、NIC 或专用设备在路径上合并规约输入：一个数据块的输入
沿聚合树的上行分支汇合，中间结果继续向根推进，response 再返回 responder
或 requester group。聚合点之后传输的是中间结果，链路字节随之减少、规约
路径随之缩短。对上面的 ring 例子，在网聚合把每个 rank 的线量从
$2(N-1)M/N$ 降回 $M$，吞吐上界从 $N/(2(N-1))$ 倍线速回到线速；
\cref{fig:baseline-anchor} 的 200-Gbps 单层控制中，NCCL-like 与三个
在网聚合模型分别达到 103.73 和 196.25--198.43~Gbps。

\section{专用聚合状态不能充分转化为可用容量}
\label{ssec:inc-survey}

数据面为每轮聚合维护上下文，通常包括通信标识、packet offset、已到达
贡献数和部分结果。实现把上下文放入一组 slot，交换机 SRAM、NIC 片上存储
或专用加速器内存决定可以同时驻留多少轮聚合；一个上下文从取得 slot 到本级
聚合完成的时间构成复用保护区间。已有系统把这些状态放在交换机 SRAM、
NIC-FPGA 或专用 reduction hardware 中~\cite{switchml,atp,netreduce,sharp,epic}。
问题在于，增加 FPGA 或片上 SRAM 扩大的只是物理容量，某个任务或路径预留的
空闲地址并不会自动服务其它活跃任务。部署成本因而同时取决于 slot 总数、
复用是否安全，以及活跃任务能否寻址已经配置的容量。

\parab{安全深度由路径、生命周期和分配序列共同决定.}
同一 offset 的输入沿多条分支进入聚合树，并在每个共享子树中合并，各分支
必须在同一级使用同一个聚合位置；选树因此同时确定报文经过的链路和逐级
状态位置。对一个聚合上下文，同一 allocator 在其复用保护区间内发生的
allocation 数决定循环映射所需的最小深度。多个通信组共享 allocator 时，
它们的 allocation 合并为同一个推进序列，安全深度由这条合并序列与上下文
生命周期共同决定；达到该深度之后，地址所有权和选择方式进一步决定有多少
slot 真正可被活跃任务使用。

\parab{现有方案在隔离与共享之间取舍.}
SwitchML~\cite{switchml}、NetReduce~\cite{netreduce}、SHARP~\cite{sharp}
和 EPIC~\cite{epic} 在预先建立的任务、连接、路径或通信组上下文中复用
状态。隔离让不同上下文独立推进，代价是静默上下文里的地址借不给活跃任务
——SwitchML 的任务分区还要同时承载 active accumulator 与 retained
result。ATP~\cite{atp} 转而允许各任务哈希寻址共享池，命中时在数据面
聚合，冲突时交由参数服务器补全。它先按 job/sequence 计算候选位置，只有
候选位置发生冲突后才执行 remap/fallback，而不是在每次 allocation 时直接
取得池中的空闲位置；共享池未满时因而仍可能碰撞。A2TP~\cite{a2tp} 和
PA-ATP~\cite{paatp} 在此基础上加入聚合器感知的公平性与进度调度。
AIRP~\cite{airp} 则由控制面联合规划拓扑、放置和聚合器内存，用全局视图
换来共享，代价是较粗的分配粒度和更新周期。

\begin{figure}[t]
\centering
\begin{tikzpicture}[x=0.25cm,y=0.72cm,font=\scriptsize]
% All three rows contain the same sixteen physical slots.  Filled cells in
% the bottom two rows illustrate allocation, not instantaneous occupancy.
\node[anchor=east] at (-0.5,3.4) {Static partition};
\foreach \i in {0,...,3}  {\fill[ArborBlue!65]   (\i,3.0) rectangle +(1,0.8);}
\foreach \i in {4,...,7}  {\fill[ArborOrange!70] (\i,3.0) rectangle +(1,0.8);}
\foreach \i in {8,...,15} {\fill[ArborLightGray] (\i,3.0) rectangle +(1,0.8);}
\foreach \i in {0,...,16} {\draw[black!65,line width=0.35pt] (\i,3.0) -- (\i,3.8);}
\draw[black!65,line width=0.45pt] (0,3.0) rectangle (16,3.8);
\foreach \i in {4,8,12} {\draw[black!80,line width=0.8pt] (\i,3.0) -- (\i,3.8);}
\node at (2,3.4) {$A$};
\node at (6,3.4) {$B$};
\node[text=black!55] at (10,3.4) {$C$};
\node[text=black!55] at (14,3.4) {$D$};
\node[anchor=west,align=left] at (16.6,3.4) {8/16 addressable\\by active jobs};

\node[anchor=east] at (-0.5,2.0) {Shared hash};
\foreach \i in {0,...,15} {\fill[white] (\i,1.6) rectangle +(1,0.8);}
\fill[red!12] (4,1.6) rectangle +(1,0.8);
\foreach \i in {0,...,16} {\draw[black!65,line width=0.35pt] (\i,1.6) -- (\i,2.4);}
\draw[black!65,line width=0.45pt] (0,1.6) rectangle (16,2.4);
\draw[red!75!black,line width=0.8pt] (4.18,1.75) -- (4.82,2.25);
\draw[red!75!black,line width=0.8pt] (4.82,1.75) -- (4.18,2.25);
\node[anchor=south,text=red!75!black] at (4.5,2.43) {$A_0,B_1$};
\node[anchor=west,align=left] at (16.6,2.0) {16/16 addressable;\\collision before full};

\node[anchor=east] at (-0.5,0.6) {Cyclic allocator};
\foreach \i in {0,...,15} {\fill[white] (\i,0.2) rectangle +(1,0.8);}
\foreach \i in {0,2,4} {\fill[ArborBlue!65]   (\i,0.2) rectangle +(1,0.8);}
\foreach \i in {1,3,5} {\fill[ArborOrange!70] (\i,0.2) rectangle +(1,0.8);}
\foreach \i in {0,...,16} {\draw[black!65,line width=0.35pt] (\i,0.2) -- (\i,1.0);}
\draw[black!65,line width=0.45pt] (0,0.2) rectangle (16,1.0);
\node[font=\tiny] at (0.5,0.6) {$A_0$};
\node[font=\tiny] at (1.5,0.6) {$B_0$};
\node[font=\tiny] at (2.5,0.6) {$A_1$};
\node[font=\tiny] at (3.5,0.6) {$B_1$};
\node[font=\tiny] at (4.5,0.6) {$A_2$};
\node[font=\tiny] at (5.5,0.6) {$B_2$};
\draw[->,black!65] (0,-0.05) -- node[below,font=\tiny] {merged allocation sequence} (6,-0.05);
\node[anchor=west,align=left] at (16.6,0.6) {16/16 allocatable;\\safe reuse after $D$};
\end{tikzpicture}

\caption{同一份 16-slot 物理容量在三种分配方式下的状态示意。配置四个
job，仅 $A/B$ 活跃：静态分区把 $C/D$ 的预留（灰色）留在原地；共享哈希
允许活跃任务寻址全池，但两个 offset 仍可能在其它地址空闲时碰撞；共享
循环 allocator 把全部活跃 allocation 合并为一个序列，在安全深度 $D$ 之后
覆盖复用。填色表示地址所有权或示例 allocation，不表示某一时刻的实际
占用率。}
\label{fig:motivation-state-allocation}
\end{figure}

\begin{table}[t]
    \centering
    \caption{典型在网聚合系统的状态所有权、原语覆盖与位置选择。}
    \label{tab:inc-comparison}
    \small
    \setlength{\tabcolsep}{3pt}
    \begin{tabularx}{\linewidth}{@{}p{0.11\linewidth}p{0.13\linewidth}p{0.11\linewidth}p{0.16\linewidth}Xp{0.14\linewidth}@{}}
        \toprule
        系统 & 状态位置 & 所有权范围 & 原语覆盖 & 位置选择 & 不可用时的处理 \\
        \midrule
        SwitchML~\cite{switchml} & 交换机 SRAM & 任务级分区 & 梯度 AllReduce & PSN 分区内循环复用 & 端侧窗口等待/恢复 \\
        NetReduce~\cite{netreduce} & RDMA 路径 FPGA & 路径上下文 & AllReduce & RDMA 报文拦截 & 端点 RoCE 恢复 \\
        SHARP~\cite{sharp} & 硬件 reduction tree & 通信组映射 & Reduce/AllReduce/Barrier & 预配置 tree & 通信库与硬件处理 \\
        EPIC~\cite{epic} & 交换机/加速器 SRAM & 组与连接 context & AllReduce/Reduce/Broadcast & PSN 窗口循环复用 & mode-specific ACK/NAK \\
        ATP~\cite{atp} & 交换机 SRAM & 跨任务共享 & 梯度 AllReduce & job/sequence hash & 参数服务器补全 \\
        AIRP~\cite{airp} & 网络内聚合器 & 控制面共享 & 随底层 INC & 全局资源规划 & 依赖底层 INC \\
        \bottomrule
    \end{tabularx}
\end{table}

\Cref{fig:motivation-state-allocation} 中的比例表示\emph{活跃任务可寻址的
状态}，而非当前非空 slot 的比例。任务级分区只让活跃任务使用自己的份额；
数据面哈希可以寻址全池，却不保证未满时没有碰撞；控制面规划有全局视图，
但受分配粒度和更新周期限制。\Cref{tab:inc-comparison} 把这些机制映射到
代表性系统。

\parab{共享哈希仍会留下不可用容量.}
\Cref{fig:motivation-hash-conflicts} 固定 $N=4,J=4$、端侧控制规则和链路
配置，仅改变 ATP 的物理聚合器数 $A$。哈希冲突率总体从 $A=8$ 时的
68.31\% 降至 $A=512$ 时的 1.08\%；在 $A=1024$ 的
3,440,448 次 aggregation attempt 中未观测到哈希冲突。有限测量中的零次
冲突并非普适保证：在 $8\le A\le960$ 的其余已测配置中，静态 job/sequence
映射仍会发生冲突，并把相应 packet 转入回退路径。

\begin{figure}[t]
\centering
\begin{tikzpicture}
\begin{axis}[
  arbor axis,
  width=0.96\linewidth,height=0.53\linewidth,
  xmin=8,xmax=1024,ymin=0,ymax=70,
  xtick={8,128,256,512,768,1024},
  ytick={0,10,20,30,40,50,60,70},
  xlabel={物理聚合器数 $A$},
  ylabel={哈希冲突率（\%）},
]
\addplot[ArborOrange,thick,mark=*,mark size=1.0pt]
  table[col sep=comma,x=physicalAggregators,y=atpHashCollisionPct]
    {plots/data/pgf/motivation-atp-hash-collisions.csv};
\end{axis}
\end{tikzpicture}

\caption{ATP 的物理聚合器数与哈希冲突率。测试床固定为 $N=4,J=4$、
每 job 每 rank 64-MiB AllReduce、8-KiB payload、端侧控制规则和链路配置；
曲线扫描全部 51 组物理聚合器配置 $A$。哈希冲突率为 FPGA 记录的 collision
数占全部 root aggregation attempts 的比例，每点汇总五个 fresh process。
$A=1024$ 的 3,440,448 次 attempt 中未观测到 collision。横纵轴均为线性
坐标。}
\label{fig:motivation-hash-conflicts}
\end{figure}

安全定容与地址覆盖因此是两个不同的问题：对共享池，全部 job 的 allocation
合并为一个推进序列并决定安全深度；达到该深度后，位置选择机制才决定已配置
地址是否都能成为活跃任务的可用容量。\Cref{sssec:slot-bound} 给出循环映射
的精确条件和部署界，\cref{ssec:eval-state} 再量化三种分配方式达到同一
性能目标所需的状态。

\section{现有组内协调依赖特定传输边界}
\label{ssec:group-consistency}

点对点 ACK、ECN、窗口和路径状态天然对应一个 sender--receiver 关系或本地
传输段，而同一 offset 的聚合完成取决于整个 fan-in。已有 INC 系统并非
忽略慢分支，而是各自在点对点连接之上建立了组内协调边界。
SwitchML~\cite{switchml} 把任务的 aggregation pool 同时用作发送窗口：
worker 收到聚合结果后才复用对应 slot，慢分支最终通过 pool 的自时钟限制
全组推进。ATP~\cite{atp} 在上行聚合时合并 CE，参数服务器再通过
parameter ACK 把共同信号返回各 worker；多播之后的单条返回分支仍可再次
标记 CE，但该信号只作用于相应 worker，不会在同一轮重新合并为组信号。

EPIC~\cite{epic} 把协调边界进一步移入交换机传输端点：Mode-I 在交换机
终结连接并逐跳流控，Mode-III 维护 pipe、ACK 进度和有效 PSN 范围，其
rate-sync 选项还可由首跳交换机为超前 rank 主动生成 CNP。维持这些边界
本身就有数据面代价：SwitchML 的 pool 同时充当发送窗口，窗口深度与
SRAM 分区互相牵制；ATP 需要参数服务器及其链路承接组信号；EPIC 则在
交换机中保存连接、ACK 进度和有效 PSN 范围等聚合状态之外的传输状态。
这份代价还随原语分化——EPIC 的 AllReduce、Broadcast 和 Reduce 分别
需要 ACK 改写反射、逐级 ACK 聚合和沿树 ACK 广播，非对称原语还要对齐
各 rank 的序号漂移；每支持一种原语，交换机传输端点上就多一份数据面
处理逻辑与存储（\cref{tab:inc-comparison} 的原语覆盖列）。因此，这些
机制都能约束分支偏斜，代价是把组内推进分别绑定到 aggregation pool、
参数服务器或交换机侧传输状态，并让原语覆盖成为数据面实现的函数。下面的
实验量化这些边界需要吸收的底层代价，而不是把它们简化为``对慢分支没有
响应''。

\parab{一条慢分支仍会扩大状态的保护区间.}
无论协调边界位于 pool、参数服务器还是交换机传输端点，聚合器都必须等同一
offset 的全部输入到齐：若一条分支与其它流共享瓶颈，快分支的输入先到也
不能单独提交，部分聚合值要一直保留到慢分支补齐。协调机制允许的到达偏斜
因而直接决定状态持有时间，也决定循环 allocator 在安全复用前必须跨越的
slot 窗口。

\begin{figure}[t]
\centering
\begin{minipage}[b]{0.46\linewidth}
\centering
\begin{tikzpicture}[
  x=0.61cm,y=0.56cm,font=\scriptsize,
  arrival/.style={-stealth,line width=0.7pt},
  event/.style={black!55,densely dashed,line width=0.45pt},
  state/.style={draw=ArborBlue!80,fill=ArborBlue!16,rounded corners=1.2pt,
    minimum height=0.48cm,align=center,font=\tiny},
]
\node[font=\bfseries\scriptsize,anchor=south] at (3.25,4.05)
  {(a) One branch delays fan-in};

\node[anchor=east,font=\tiny] at (0.72,3.15) {7 fast inputs};
\draw[arrival,ArborBlue] (0.86,3.15) -- (2.05,3.15);
\node[anchor=east,font=\tiny,text=ArborOrange] at (0.72,2.35)
  {1 contended input};
\draw[arrival,ArborOrange] (0.86,2.35) -- (5.25,2.35);

\draw[event] (2.05,0.55) -- (2.05,3.53);
\draw[event] (5.25,0.55) -- (5.25,3.53);
\node[anchor=south,align=center,font=\tiny] at (2.05,3.56)
  {partial aggregate\\created};
\node[anchor=south,align=center,font=\tiny] at (5.25,3.56)
  {fan-in complete};

\node[anchor=east,font=\tiny] at (1.76,1.35) {Root slot};
\node[state,anchor=west,minimum width=3.20cm] at (2.05,1.35)
  {partial result remains live};
\draw[-stealth,black!65,line width=0.45pt] (1.40,0.55) -- (5.72,0.55);
\node[anchor=north,font=\tiny] at (3.55,0.43) {Time};
\end{tikzpicture}
\end{minipage}\hfill
\begin{minipage}[b]{0.51\linewidth}
\centering
\begin{tikzpicture}
\begin{axis}[
  arbor axis,arbor panel,
  width=\linewidth,height=0.50\columnwidth,
  title={(b) Measured state expansion},
  xmin=0.8,xmax=4.6,
  ymin=-0.45,ymax=1.45,
  xtick={1,2,3,4},
  xticklabels={1x,2x,3x,4x},
  ytick={0,1},
  yticklabels={{State holding p99\\$31.8\!\rightarrow\!127.2~\mu$s},
    {Safe reuse depth\\$89\!\rightarrow\!229$ slots}},
  y dir=reverse,
  xmajorgrids=true,
  ymajorgrids=false,
  yticklabel style={font=\scriptsize,align=right},
  xlabel={Increase over no contention},
]
\addplot[black!70,thin,densely dotted,mark=none,forget plot]
  coordinates {(1,-0.45) (1,1.45)};
\addplot[
  ArborGray,only marks,mark=square*,mark size=2.0pt,forget plot,
]
  table[col sep=comma,x expr=1,y=x]
    {plots/data/pgf/motivation-branch-skew.csv};
\addplot[
  ArborBlue,only marks,mark=*,mark size=2.4pt,
  nodes near coords,
  point meta=explicit symbolic,
  every node near coord/.append style={font=\scriptsize,anchor=west,xshift=2pt},
  error bars/.cd,x dir=both,x explicit,
]
  table[col sep=comma,x=relative,y=x,x error=relativeCi,meta=label]
    {plots/data/pgf/motivation-branch-skew.csv};
\end{axis}
\end{tikzpicture}
\end{minipage}

\caption{单条受竞争分支扩大聚合状态的保护区间。(a) 同一 offset 的七个
输入先到，root slot 持续保存部分结果，直到受竞争分支的最后一个输入到达。
(b) root 状态的 p99 持有时间由 31.8 增至 127.2~\textmu s，安全循环复用
深度由 89 增至 229 slots，分别为无背景流时的 4.0 和 2.6 倍。点为三条配对
物理分支的均值，误差线为 Student-t 95\% CI。}
\label{fig:motivation-branch-skew}
\end{figure}

为隔离这一效应，\cref{fig:motivation-branch-skew} 把一个 8-rank
AllReduce 固定到同一棵树上，只加入一条持续点对点流，与其中一条 200-Gbps
上行链路竞争。root 状态的 p99 持有时间从 31.83 增至 127.24~\textmu s，
完整 slot 事件轨迹给出的安全循环复用深度从 89 增至 229；树、聚合位置和
其它链路均保持不变，8192-slot 配置不构成瓶颈，全部运行也没有丢包、PFC、
\texttt{AGG\_MISS}、overwrite 或 repair。增长因此只来自单条分支拖延整组
完成，与聚合状态不足或恢复流量无关。这个对照实验固定发放行为，没有复现
上述系统完整的自时钟、参数服务器或逐跳流控；它量化的是这些机制共同需要
约束的状态生命周期，其中的 p99 也只是有限运行中的观测值，不是定容所需的
硬上界。\Cref{ssec:eval-cc} 在相同的固定单树、禁止换路条件下，进一步
比较不同反馈范围如何控制发送。

\parab{候选树之间的剩余容量不一致.}
ECMP~\cite{ecmp} 在 flow 粒度固定路径；Themis~\cite{themis} 用确定性的
PSN--path 映射喷洒 packet，常规分配同样由报文标识决定，不响应持续的路径
不对称。CONGA~\cite{conga} 和 LetFlow~\cite{letflow} 等自适应方案按网络
反馈在 flowlet 或 packet 粒度迁移后续流量，但一次决策只移动一个单播
packet 或 flowlet。聚合树的剩余容量不同时，固定映射会把一部分 offset 留
在慢树上；自适应单播能迁移后续流量，却决定不了同一 offset 的整组输入该
使用哪棵聚合树。

\parab{在网 fan-in 改变选路的决策单位与一致性义务.}
单播多路径中，一次选路决策的错误代价是乱序或次优路径
（\cref{ssec:p2p-ccl}）。在网 fan-in 下，为一个 packet offset 选择聚合树
是一次整组决策：它同时确定整组输入经过的链路和逐级聚合位置。这一决定必须
传达给全部 requester，而不同分支上的控制信息可能以不同顺序到达。若各
requester 仅按本地连接状态把一次选树决定绑定到``下一个未发送 packet''，
同一决定就可能对应不同 offset，使一轮 fan-in 合并不同数据的输入。一棵
固定树天然避免该错误，却无法利用其它树的动态剩余容量；静态条带化也只能
在预定比例内保持一致。本文把自适应方案必须维持的关系称为\emph{聚合上下文
一致性}：同一 packet offset 的全部输入必须导出同一棵树和同一份逐级聚合
上下文。

\section{上下文与反馈要求及多播-聚合流}
\label{ssec:inc-limits}

传输层因此需要满足两项要求。\emph{聚合上下文一致性}要求同一 offset 的
输入进入同一棵树，并在共享的每一级使用一致的聚合位置；\emph{整组进度
可见性}要求发送控制、后续选树和恢复依据覆盖最慢分支的完成边界或组级拥塞
信号。多租户状态管理还需独立保证上下文在复用前已经安全结束，动态选树才
不必把``是否还有空闲 slot''当作又一个在线决策变量。

本文把同时描述一对多发送与多对一聚合的传输关系称为\emph{多播-聚合流}。
它固定一个 responder 和 requester group，并以 packet offset 作为路径、
聚合上下文、组内推进和恢复的共同语义边界；\cref{chap:design} 给出具体
实现。
